% SC26 MLIR tutorial -- Session 1, Part 2 (GPU). Presenter deck.
%
% Build:  make            (the deck PDF)
% See presentation/README.md.
\providecommand{\notesmode}{0}
\documentclass[aspectratio=169,10pt]{beamer}
\usepackage{sc26}
\ifnum\notesmode=1 \setbeameroption{show notes}\fi

\title{Introduction to MLIR}
\subtitle{Session 1, Part 2 --- GPU}
\author{SC26 Tutorial}
\date{}

\begin{document}

\begin{frame}[plain]
  \titlepage
\end{frame}
\note{Pre-run setup cells 0--7 of the walkthrough notebook on a T4 runtime
before the room fills: the toolchain download and the first plugin build are
the slowest things in this part and neither teaches anything.

One exercise in this part, the GPU lowering: about 10 minutes of keyboard
time. Step 1 is a kernel you walk through; nobody types anything in it.}

% =======================================================================
\begin{frame}{Setup}
  \begin{itemize}
    \item<1-> \textbf{Switch the runtime to a GPU:}
      \cd{Change runtime type $\to$ T4 GPU}, before running any cell.
    \item<2-> \textbf{Fresh runtime.} Nothing from Part 1 carries
      over, so the notebook downloads the toolchain and builds the plugin
      again.
    \item<3-> \textbf{We start from the Part 1 solution.} The notebook clones the
      \cd{session1-part1-solution} branch, where \cd{blas.matmul} is already
      defined. Nothing to re-type.
    \item<4-> \textbf{One exercise:} GPU lowering.
  \end{itemize}
\end{frame}
\note{Change runtime type is under the Runtime menu, or the arrow next to
Connect at the top right. Ctrl+F8 (run every cell above) still works if
someone loses their place.}

% =======================================================================
\begin{frame}{GPU Lowering}
  \begin{center}
  \begin{tikzpicture}[
      box/.style={draw=scblue,fill=scwash,rounded corners=2pt,inner sep=5pt,
                  font=\ttfamily\scriptsize,minimum height=2.4em,align=center},
      lbl/.style={font=\scriptsize\itshape,color=scgrey,align=center},
      arr/.style={-{Stealth[length=2mm]},draw=scgrey,thick}]

    \node[box,draw=scaccent] (op) at (0,0) {blas.matmul};
    \node[box] (scf) at (4.6,1.0) {scf.for nest};
    \node[box] (cpu) at (10.6,1.0) {JIT + run\\\normalfont\itshape CPU};
    \draw[arr] (op) -- node[lbl,above left=-2pt and -8pt]
      {convert-blas-to-scf\\Part 1} (scf);
    \draw[arr] (scf) -- node[lbl,above]{7 passes} (cpu);

    \visible<2->{
      \node[box,draw=scaccent] (gpu) at (4.6,-1.0) {gpu.launch kernel};
      \node[box] (t4) at (10.6,-1.0) {JIT + run\\\normalfont\itshape T4 GPU};
      \draw[arr,draw=scaccent] (op) -- node[lbl,below left=-2pt and -8pt]
        {convert-blas-to-gpu\\\emph{this part}} (gpu);
      \draw[arr] (gpu) -- node[lbl,below]{gpu-lower-to-nvvm-pipeline} (t4);
    }
  \end{tikzpicture}
  \end{center}
  \vspace{0.3em}
  \begin{itemize}
    \item<1-> In Part 1, we lowered \cd{blas.matmul} to loops and ran it on
      the CPU.
    \item<2-> In this part, we lower the \alert{same op} to a \cd{gpu.launch}
      kernel and run it on a GPU.
    \item<3-> First we read that kernel written by hand, so you know what
      your pattern has to build.
  \end{itemize}
\end{frame}
\note{The op definition and its verifier do not change at all. Everything in
this part is a second pattern in a second pass.}

% =======================================================================
\section{Step 1 --- Loop-Level GPU Kernel}
\begin{frame}[plain]\sectionpage\end{frame}

\begin{frame}{Blocks and Threads}
  \begin{center}
  \begin{tikzpicture}[
      blk/.style={draw=scblue,fill=scwash,minimum width=1.25cm,
                  minimum height=0.8cm,font=\ttfamily\scriptsize},
      thr/.style={draw=scaccent,fill=white,minimum size=0.42cm,inner sep=0pt},
      lbl/.style={font=\scriptsize,color=scink}]

    \foreach \x in {0,1,2} {
      \foreach \y in {0,1} {
        \node[blk] (b\x\y) at (\x*1.4, -\y*0.95) {(\y, \x)};
      }
    }
    \node[lbl,above=0.15cm of b10] {a grid of blocks --- \cd{blockIdx} as (y, x)};

    \visible<2->{
      \foreach \x in {0,...,4} {
        \foreach \y in {0,...,3} {
          \node[thr] (t\x\y) at (6.4+\x*0.48, 0.25-\y*0.48) {};
        }
      }
      \node[lbl,above=0.15cm of t20] {one block of threads --- \cd{threadIdx}};
      \draw[scgrey,dashed] (b20.north east) -- (t00.north west);
      \draw[scgrey,dashed] (b20.south east) -- (t03.south west);
      \draw[-{Stealth[length=1.5mm]},scgrey] (6.15,-1.6) -- (8.45,-1.6)
        node[below,font=\scriptsize]{x};
      \draw[-{Stealth[length=1.5mm]},scgrey] (8.8,0.45) -- (8.8,-1.35)
        node[right,font=\scriptsize]{y};
    }
  \end{tikzpicture}
  \end{center}
  \vspace{0.4em}
  \begin{itemize}
    \item<1-> A kernel runs the \alert{same body} on many threads at once.
    \item<2-> Threads are grouped into \alert{blocks}, and blocks into a
      \alert{grid}. Each can be up to 3D: x, y and z.
    \item<3-> Every thread knows its \cd{blockIdx}, its \cd{threadIdx} and the
      block size, \cd{blockDim}. Its global position in each dimension is
      \cd{blockIdx * blockDim + threadIdx}.
  \end{itemize}
\end{frame}
\note{This is the CUDA model, and the gpu dialect uses it as is. Our kernel
is the simplest case: one block, so blockIdx is always zero and threadIdx
alone says which element a thread computes.}

% =======================================================================
\begin{frame}[fragile]{\cd{gpu.launch}}
\begin{lstlisting}[language=mlir]
(@\lkw{gpu.launch} \hl<4>{blocks(\%bx, \%by, \%bz)} \hl<2>{in (\%gridX = \%c1, \%gridY = \%c1, \%gridZ = \%c1)}@)
           (@\hl<4>{threads(\%tx, \%ty, \%tz)} \hl<3>{in (\%blockX = \%c5, \%blockY = \%c4, \%blockZ = \%c1)} \{@)
  (@\hl<5>{\lcm{// the body: run once by every thread}}@)
  (@\hl<5>{\lkw{gpu.terminator}}@)
}
\end{lstlisting}
  \vspace{0.4em}
  \begin{itemize}
    \item<2-> \cd{blocks(...) in (...)}: how many blocks, in x, y and z. Here,
      one.
    \item<3-> \cd{threads(...) in (...)}: how many threads in each block. Here
      5 $\times$ 4 = 20. Inside the body \cd{\%blockX, \%blockY, \%blockZ}
      name those sizes: \cd{blockDim}.
    \item<4-> \cd{\%bx, \%by, \%bz} and \cd{\%tx, \%ty, \%tz} are bound inside
      the body: the block and thread ids, CUDA's \cd{blockIdx} and
      \cd{threadIdx}.
    \item<5-> Every thread runs the body, and \cd{gpu.terminator} ends it.
  \end{itemize}
\end{frame}
\note{The sizes are ordinary index values, so a pattern can compute them --
which is exactly what the one in step 2 does, from the shapes of the
operands.}

% =======================================================================
\begin{frame}[fragile]{\cd{matmul\_gpu.mlir}, Step 1 of 2 --- Launch}
\begin{lstlisting}[language=mlir]
func.func @matmul(%a: memref<4x3xf32>, %b: memref<3x5xf32>,
                  %out: memref<4x5xf32>) {
  (@\hl<2>{\%a\_u = \lkw{memref.cast} \%a : \lkw{memref}<4x3xf32> \lkw{to} \lkw{memref}<*xf32>}@)
  (@\hl<2>{\%b\_u = \lkw{memref.cast} \%b : \lkw{memref}<3x5xf32> \lkw{to} \lkw{memref}<*xf32>}@)
  (@\hl<2>{\%out\_u = \lkw{memref.cast} \%out : \lkw{memref}<4x5xf32> \lkw{to} \lkw{memref}<*xf32>}@)
  (@\hl<2>{\lkw{gpu.host\_register} \%a\_u : \lkw{memref}<*xf32>}@)
  (@\hl<2>{\lkw{gpu.host\_register} \%b\_u : \lkw{memref}<*xf32>}@)
  (@\hl<2>{\lkw{gpu.host\_register} \%out\_u : \lkw{memref}<*xf32>}@)
  %c1 = arith.constant 1 : index
  %c3 = arith.constant 3 : index
  %c4 = arith.constant 4 : index
  %c5 = arith.constant 5 : index
  (@\hl<3>{\lkw{gpu.launch} blocks(\%bx, \%by, \%bz) in (\%gridX = \%c1, \%gridY = \%c1, \%gridZ = \%c1)}@)
             (@\hl<3>{threads(\%tx, \%ty, \%tz) in (\%blockX = \%c5, \%blockY = \%c4, \%blockZ = \%c1) \{}@)
    // the body: next slides
    gpu.terminator
  }
  return
}
\end{lstlisting}
  \vspace{0.2em}
  \begin{itemize}
    \item<2-> The \cd{memref.cast} + \cd{gpu.host\_register} lines make the
      host buffers reachable from the GPU.
    \item<3-> One block of 5 $\times$ 4 threads: one thread per element of the
      4 $\times$ 5 \cd{out}.
  \end{itemize}
\end{frame}
\note{session1/matmul\_gpu.mlir ships complete. The cast to an unranked memref
is only because gpu.host\_register takes one. If someone asks what happens
without the registration: on some GPU and driver setups the kernel silently
reads and writes garbage rather than faulting.}

% =======================================================================
\begin{frame}{Thread to Element Mapping}
  \begin{center}
  \begin{tikzpicture}[
      cell/.style={draw=scblue,fill=scwash,minimum width=1.15cm,
                   minimum height=0.55cm,font=\ttfamily\scriptsize},
      hdr/.style={font=\ttfamily\scriptsize,color=scaccent}]
    \foreach \c in {0,...,4} {
      \node[hdr] at (\c*1.25, 0.55) {tx = \c};
      \foreach \r in {0,...,3} {
        \node[cell] at (\c*1.25, -\r*0.65) {(\r, \c)};
      }
    }
    \foreach \r in {0,...,3} {
      \node[hdr,anchor=east] at (-0.75, -\r*0.65) {ty = \r};
    }
    \node[font=\scriptsize,color=scgrey,anchor=west] at (6.0, -0.98)
      {\cd{out}, 4 $\times$ 5: each cell is};
    \node[font=\scriptsize,color=scgrey,anchor=west] at (6.0, -1.33)
      {the thread \cd{(ty, tx)} that computes it};
  \end{tikzpicture}
  \end{center}
  \vspace{0.5em}
  \begin{itemize}
    \item<2-> y spans the 4 rows of \cd{out}, x its 5 columns. Each thread
      computes one element of the output.
  \end{itemize}
\end{frame}
\note{The coalescing argument, without timing (Colab timings are noise):
memrefs are row-major, so out[i, j] and b[k, j] are contiguous in j. A warp
is 32 threads with consecutive threadIdx.x. With x on the columns those lanes
touch adjacent floats, one memory transaction; with x on the rows each lane
is a whole row away, one transaction per lane. Session 2 Part 2 comes back to
this when a pass picks the mapping instead.}

% =======================================================================
\begin{frame}[fragile]{\cd{matmul\_gpu.mlir}, Step 2 of 2 --- Body}
\begin{lstlisting}[language=mlir,basicstyle=\ttfamily\fontsize{6.4}{7.6}\selectfont]
gpu.launch blocks(%bx, %by, %bz) ... threads(%tx, %ty, %tz) ... {
  (@\hl<2>{\%rowBase = \lkw{arith.muli} \%by, \%blockY : \lkw{index}}@)
  (@\hl<2>{\%row = \lkw{arith.addi} \%rowBase, \%ty : \lkw{index}}@)
  (@\hl<2>{\%colBase = \lkw{arith.muli} \%bx, \%blockX : \lkw{index}}@)
  (@\hl<2>{\%col = \lkw{arith.addi} \%colBase, \%tx : \lkw{index}}@)
  %c0 = arith.constant 0 : index
  %step = arith.constant 1 : index
  (@\hl<3>{\%init = \lkw{memref.load} \%out[\%row, \%col] : \lkw{memref}<4x5xf32>}@)
  (@\hl<4>{\%acc = \lkw{scf.for} \%k = \%c0 \lkw{to} \%c3 \lkw{step} \%step \lkw{iter\_args}(\%s = \%init) -> (\lkw{f32}) \{}@)
    (@\hl<4>{\%av = \lkw{memref.load} \%a[\%row, \%k] : \lkw{memref}<4x3xf32>}@)
    (@\hl<4>{\%bv = \lkw{memref.load} \%b[\%k, \%col] : \lkw{memref}<3x5xf32>}@)
    %prod = arith.mulf %av, %bv : f32
    %s2 = arith.addf %s, %prod : f32
    scf.yield %s2 : f32
  }
  (@\hl<3>{\lkw{memref.store} \%acc, \%out[\%row, \%col] : \lkw{memref}<4x5xf32>}@)
  gpu.terminator
}
\end{lstlisting}
  \vspace{0.2em}
  \begin{itemize}
    \item<2-> Global row and column: block id $\times$ block size + thread id.
    \item<3-> Thread \cd{(\%ty, \%tx)} reads, accumulates and writes
      \cd{out[\%row, \%col]}.
    \item<4-> The reduction over \cd{\%k} is Part 1's loop, on row \cd{\%row}
      of \cd{a} and column \cd{\%col} of \cd{b}.
    \item<5-> No explicit loops in the code over rows and columns:
      \alert{the 2D grid of threads spans those loops}.
  \end{itemize}
\end{frame}
\note{The block-id terms are zero here, so this computes the same as indexing
with \%ty and \%tx directly. It is written out because it is the general
form: the same body is right for any number of blocks.}

% =======================================================================
\begin{frame}[fragile]{Lowering and Execution}
\begin{lstlisting}[style=plain,basicstyle=\ttfamily\small]
gpu-lower-to-nvvm-pipeline{cubin-chip=sm_75 cubin-format=bin
                           cubin-features=+ptx78}
\end{lstlisting}
  \vspace{0.4em}
  \begin{itemize}
    \item<2-> Moves the launch body out into a separate kernel, lowers it to
      NVVM, and compiles it to a binary for this GPU.
    \item<3-> Rewrites the host side to launch that kernel through the CUDA
      driver.
    \item<4-> The JIT needs the CUDA runtime wrapper library in
      \cd{shared\_libs}, and \cd{engine.initialize()} to load the compiled
      kernel.
  \end{itemize}
  \vspace{0.4em}
  \visible<5->{\begin{center}\color{scgreen}\large\ttfamily OK\end{center}}
\end{frame}
\note{The notebook reads the chip from nvidia-smi rather than hard-coding
sm\_75, so it also works on whatever GPU Colab hands out.

Without engine.initialize() every CUDA call fails with
CUDA\_ERROR\_INVALID\_HANDLE: the compiled kernel is loaded by a global
constructor that nothing else runs.}

% =======================================================================
\section{Step 2 --- Lowering \texttt{blas.matmul} to GPU}
\begin{frame}[plain]\sectionpage\end{frame}

\begin{frame}[fragile]{Reference: \cd{MatVecToGPUPattern}, Step 1 of 5 --- Setup}
  \begin{columns}[t]
    \begin{column}{0.59\textwidth}
      \begin{center}\small\textbf{the C++}\end{center}
\begin{lstlisting}[style=halfs,language=C++,morekeywords={LogicalResult,Value,Location,override,struct,public,using}]
struct MatVecToGPUPattern : public OpRewritePattern<MatVecOp> {
  using OpRewritePattern::OpRewritePattern;
  LogicalResult
  matchAndRewrite(MatVecOp op,
                  PatternRewriter &rewriter) const override {
    Location loc = op.getLoc();
    Value mat = op.getMat();
    Value vec = op.getVec();
    Value out = op.getOut();
    (@\hl<2>{\lcm{// memref.cast + gpu::HostRegisterOp for mat, vec, out}}@)
    Value c0 = arith::ConstantIndexOp::create(rewriter, loc, 0);
    Value c1 = arith::ConstantIndexOp::create(rewriter, loc, 1);
    Value numRows = memref::DimOp::create(rewriter, loc, mat, c0);
    Value numCols = memref::DimOp::create(rewriter, loc, mat, c1);
    // steps 2-4 go here
    rewriter.eraseOp(op);
    return success();
  }
};
\end{lstlisting}
    \end{column}
    \begin{column}{0.40\textwidth}
      \begin{center}\small\textbf{the IR it generates}\end{center}
\begin{lstlisting}[style=halfs,language=mlir]
(@\strut@)








(@\hl<2>{\lcm{// memref.cast + gpu.host\_register}}@)
%c0 = arith.constant 0 : index
%c1 = arith.constant 1 : index
%numRows = memref.dim %mat, %c0
%numCols = memref.dim %mat, %c1
\end{lstlisting}
    \end{column}
  \end{columns}
  \vspace{0.3em}
  \begin{itemize}
    \item<2-> \cd{MatVecToSCFPattern}'s start, plus step 1's host
      registration (six lines, folded here).
    \item<3-> In \cd{MatMulToGPUPattern} all of this is \alert{provided}. You
      start below it.
  \end{itemize}
\end{frame}
\note{\cd{blas-dialect/lib/LowerBlasToGPU.cpp}. MatVecToGPUPattern is complete
in the file. As in Part 1, the right-hand side shows what the C++ builds; the
pass's printed output folds the memref.dim ops to constants.}

% =======================================================================
\begin{frame}[fragile]{Reference: \cd{MatVecToGPUPattern}, Step 2 of 5 --- Launch}
  {\small\textbf{the C++}}
\begin{lstlisting}[language=C++,basicstyle=\ttfamily\fontsize{6}{7.2}\selectfont,morekeywords={auto,Value}]
auto launchOp =
    (@\hl<2>{gpu::LaunchOp::create(rewriter, loc, c1, c1, c1, numRows, c1, c1);}@)
{
  OpBuilder::InsertionGuard guard(rewriter);
  (@\hl<3>{rewriter.setInsertionPointToStart(\&launchOp.getBody().front());}@)
  (@\hl<4>{\lkw{Value} bx = launchOp.getBlockIds().x;}@)
  (@\hl<4>{\lkw{Value} blockDimX = launchOp.getBlockSize().x;}@)
  (@\hl<4>{\lkw{Value} tx = launchOp.getThreadIds().x;}@)
  (@\hl<5>{\lkw{Value} rowBase = arith::MulIOp::create(rewriter, loc, bx, blockDimX);}@)
  (@\hl<5>{\lkw{Value} row = arith::AddIOp::create(rewriter, loc, rowBase, tx);}@)
  // steps 3 and 4 go here
}
\end{lstlisting}
  {\small\textbf{the IR it generates}}
\begin{lstlisting}[language=mlir,basicstyle=\ttfamily\fontsize{6}{7.2}\selectfont]
(@\hl<2>{\lkw{gpu.launch} blocks(}\hl<2,4>{\%bx}\hl<2>{, \%by, \%bz) in (\%gridX = \%c1, \%gridY = \%c1, \%gridZ = \%c1)}@)
           (@\hl<2>{threads(}\hl<2,4>{\%tx}\hl<2>{, \%ty, \%tz) in (}\hl<2,4>{\%blockX}\hl<2>{ = \%numRows, \%blockY = \%c1, \%blockZ = \%c1) \{}@)
  (@\hl<5>{\%rowBase = \lkw{arith.muli} \%bx, \%blockX : \lkw{index}}@)
  (@\hl<5>{\%row = \lkw{arith.addi} \%rowBase, \%tx : \lkw{index}}@)
  (@\hl<3>{\lcm{// steps 3 and 4 go here}}@)
}
\end{lstlisting}
  \vspace{0.1em}
  \begin{itemize}
    \item<2-> Grid size: \cd{(1, 1, 1)}, block size: \cd{(numRows, 1, 1)}.
    \item<3-> The body is a region: move the insertion point into it, as with
      Part 1's loops.
    \item<4-> Then the block id, block size and thread id: \cd{bx},
      \cd{blockDimX}, \cd{tx}.
    \item<5-> From them the global row, \cd{bx * blockDim.x + tx}.
  \end{itemize}
\end{frame}
\note{One thread per row of mat: the threads take the place of Part 1's row
loop. Writing the row as block id times block size plus thread id is the
canonical form; it is what makes the same pattern right when the launch uses
more than one block.}

% =======================================================================
\begin{frame}[fragile]{Reference: \cd{MatVecToGPUPattern}, Step 3 of 5 --- Body}
  \begin{columns}[t]
    \begin{column}{0.59\textwidth}
      \begin{center}\small\textbf{the C++}\end{center}
\begin{lstlisting}[style=halfs,language=C++,morekeywords={auto,Value}]
Value init = memref::LoadOp::create(rewriter, loc, out,
                                    ValueRange{(@\df{row}@)});
auto redLoop = scf::ForOp::create(rewriter, loc, c0, numCols, c1,
                                  ValueRange{init});
{
  // the guard, the insertion point, j and acc: as in Part 1
  Value matElem = memref::LoadOp::create(rewriter, loc, mat,
                                         ValueRange{(@\df{row}@), j});
  Value vecElem = memref::LoadOp::create(rewriter, loc, vec,
                                         ValueRange{j});
  // prod, sum and the yield: as in Part 1
}
\end{lstlisting}
    \end{column}
    \begin{column}{0.40\textwidth}
      \begin{center}\small\textbf{the IR it generates}\end{center}
\begin{lstlisting}[style=halfs,language=mlir]
%init = memref.load %out[(@\df{\%row}@)]

%acc = scf.for %j = %c0 to %numCols step %c1
    iter_args(%a = %init) -> (f32) {

  ...
  %m = memref.load %mat[(@\df{\%row}@), %j]

  %v = memref.load %vec[%j]

  ...
}
\end{lstlisting}
    \end{column}
  \end{columns}
  \vspace{0.6em}
  \begin{itemize}
    \item<2-> This is \cd{MatVecToSCFPattern}'s body from Part 1, with the row
      index replaced: \cd{row} where Part 1 had \cd{i}.
    \item<3-> There is no row loop: \alert{the threads span the instances of the
      row loop}.
  \end{itemize}
\end{frame}

% =======================================================================
\begin{frame}[fragile]{Reference: \cd{MatVecToGPUPattern}, Step 4 of 5 --- Finish the Body}
  \begin{columns}[t]
    \begin{column}{0.62\textwidth}
      \begin{center}\small\textbf{the C++}\end{center}
\begin{lstlisting}[style=halfs,language=C++,morekeywords={LogicalResult,auto,Value,override}]
LogicalResult matchAndRewrite(MatVecOp op, ...) const override {
  // step 1: operands, host registration, c0, c1, numRows, numCols
  auto launchOp = gpu::LaunchOp::create(..., numRows, c1, c1);
  {
    // step 2: the insertion point, row = bx * blockDim.x + tx
    Value init = memref::LoadOp::create(..., out, ValueRange{row});
    auto redLoop = scf::ForOp::create(..., ValueRange{init});
    {
      // step 3: the reduction body
    }
    (@\df{memref::StoreOp::create(rewriter, loc, redLoop.getResult(0), out,}@)
                            (@\df{ValueRange\{row\});}@)
    (@\df{gpu::TerminatorOp::create(rewriter, loc);}@)
  }
  rewriter.eraseOp(op);
  return success();
}
\end{lstlisting}
    \end{column}
    \begin{column}{0.37\textwidth}
      \begin{center}\small\textbf{the IR it generates}\end{center}
\begin{lstlisting}[style=halfs,language=mlir]
(@\strut@)

gpu.launch ... {


  %init = memref.load %out[%row]
  %acc = scf.for %j ... {

    ...
  }
  (@\df{memref.store \%acc, \%out[\%row]}@)

  (@\df{gpu.terminator}@)
}
\end{lstlisting}
    \end{column}
  \end{columns}
  \vspace{0.5em}
  \begin{itemize}
    \item<2-> The store is Part 1's, indexed by \cd{row}: after the reduction
      loop, inside the launch body.
    \item<3-> A launch body must end in \cd{gpu.terminator}, and nothing
      creates it for you.
  \end{itemize}
\end{frame}
\note{Unlike an scf.for without iter\_args, gpu::LaunchOp::create leaves the
body without a terminator. Forget it and the verifier rejects the launch.}

% =======================================================================
\begin{frame}[fragile]{Reference: \cd{MatVecToGPUPattern}, Step 5 of 5 --- Registration}
  \cd{LowerBlasToGPU.cpp}, the GPU pass:
\begin{lstlisting}[language=C++,basicstyle=\ttfamily\scriptsize,morekeywords={void,override}]
void runOnOperation() override {
  RewritePatternSet patterns(&getContext());
  (@\hl<2>{patterns.add<MatVecToGPUPattern>(\&getContext());}@)
  // TODO(exercise): register your pattern too:
  //   patterns.add<MatMulToGPUPattern>(&getContext());
  ...
}
\end{lstlisting}
  \vspace{0.3em}
  \begin{itemize}
    \item<2-> Registered exactly as in Part 1. Yours goes on the same line.
  \end{itemize}
\end{frame}

% =======================================================================
\begin{frame}{Exercise 5 --- Write the GPU Lowering for Matmul}
  \begin{center}\large
    \cd{MatMulToGPUPattern} in \cd{blas-dialect/lib/LowerBlasToGPU.cpp}
  \end{center}
  \vspace{0.3em}
  \begin{itemize}
    \item \textbf{Provided}, commented out: the operands, the host
      registration, \cd{c0}, \cd{c1}, \cd{M}, \cd{K}, \cd{N}, and the
      \cd{eraseOp}. Uncomment it.
    \item \textbf{You write} the \cd{gpu::LaunchOp::create} and its body ---
      the load, the reduction loop, the store, the \cd{gpu::TerminatorOp} ---
      then register the pattern.
  \end{itemize}
  \vspace{0.3em}
  \textbf{What is different from matvec:}
  \begin{itemize}
    \item A 2D thread grid: \cd{N} threads in x (the columns of \cd{out}),
      \cd{M} in y (its rows) --- the grid from step 1.
    \item Two global indices: \cd{row} from the y ids and \cd{col} from the x
      ids, each built the way matvec builds its \cd{row}.
    \item The reduction runs over \cd{K}, reading \cd{a[row, k]} and
      \cd{b[k, col]}, into \cd{out[row, col]}.
  \end{itemize}
  \vspace{0.4em}
  \begin{center}\color{scaccent}\large
    10 minutes to write it --- then we walk through the solution.
  \end{center}
\end{frame}
\note{If you are behind schedule, type this one yourself and narrate.

The check is the rebuild cell, then the cell that runs the pass and the one
that runs its output on the GPU.}

% =======================================================================
\begin{frame}[fragile]{Solution --- \cd{MatMulToGPUPattern}, and What Changed from Matvec}
\begin{lstlisting}[language=C++,morekeywords={auto,Value},basicstyle=\ttfamily\fontsize{6.6}{7.9}\selectfont]
auto launchOp =        (@\only<2>{\color{scaccent}$\leftarrow$ a 2D grid: N threads in x, M in y}@)
    gpu::LaunchOp::create(rewriter, loc, c1, c1, c1, (@\df{N, M}@), c1);
{
  OpBuilder::InsertionGuard guard(rewriter);
  rewriter.setInsertionPointToStart(&launchOp.getBody().front());
  // by, bx, blockDimY, blockDimX, ty, tx: as matvec gets bx, blockDimX, tx
  Value rowBase = arith::MulIOp::create(rewriter, loc, (@\df{by, blockDimY}@));    (@\only<3>{\color{scaccent}$\leftarrow$ row from y}@)
  Value row = arith::AddIOp::create(rewriter, loc, rowBase, (@\df{ty}@));
  (@\df{Value colBase = arith::MulIOp::create(rewriter, loc, bx, blockDimX);}@)  (@\only<3>{\color{scaccent}$\leftarrow$ col from x}@)
  (@\df{Value col = arith::AddIOp::create(rewriter, loc, colBase, tx);}@)

  Value init =
      memref::LoadOp::create(rewriter, loc, out, ValueRange{row(@\df{, col}@)});
  auto redLoop =
      scf::ForOp::create(rewriter, loc, c0, (@\df{K}@), c1, ValueRange{init});
  {
    // the guard, the insertion point, k and acc: as in matvec
    Value aElem = memref::LoadOp::create(rewriter, loc, a, ValueRange{row, k});
    Value bElem = memref::LoadOp::create(rewriter, loc, b, ValueRange{k(@\df{, col}@)});
    // prod, sum and the yield: as in matvec
  }
  memref::StoreOp::create(rewriter, loc, redLoop.getResult(0), out,
                          ValueRange{row(@\df{, col}@)});     (@\only<4>{\color{scaccent}$\leftarrow$ two indices}@)
  gpu::TerminatorOp::create(rewriter, loc);
}
\end{lstlisting}
\end{frame}
\note{This goes where the scaffolding says "your gpu.launch here". Everything
in black is the matvec pattern with the names changed; the orange is what is
new. It generates the step 1 kernel: one block of N by M threads, and thread
(ty, tx) computing out[row, col].}

\begin{frame}[fragile]{Solution --- Register the Pattern}
\begin{lstlisting}[language=C++,basicstyle=\ttfamily\scriptsize,morekeywords={void,override}]
void runOnOperation() override {
  RewritePatternSet patterns(&getContext());
  patterns.add<MatVecToGPUPattern(@\df{, MatMulToGPUPattern}@)>(&getContext());
  ...
}
\end{lstlisting}
  \vspace{0.25em}
  \begin{itemize}
    \item<2-> Forget it and the pass stops with:
  \end{itemize}
  \vspace{0.3em}
  \visible<2->{\errbox{error: convert-blas-to-gpu did not convert 'blas.matmul' -{}- is its pattern\\
    written and registered in blas-dialect/lib/LowerBlasToGPU.cpp?}}
\end{frame}

% =======================================================================
\begin{frame}{Recap}
  \begin{center}
  \begin{tikzpicture}[
      box/.style={draw=scblue,fill=scwash,rounded corners=2pt,inner sep=5pt,
                  font=\ttfamily\scriptsize,minimum height=2.4em,align=center},
      lbl/.style={font=\scriptsize\itshape,color=scgrey,align=center},
      arr/.style={-{Stealth[length=2mm]},draw=scgrey,thick}]

    \node[box,draw=scaccent] (op) at (0,0) {blas.matmul};
    \node[box] (scf) at (4.6,0.9) {scf.for nest};
    \node[box] (cpu) at (10.6,0.9) {JIT + run\\\normalfont\itshape CPU};
    \node[box,draw=scaccent] (gpu) at (4.6,-0.9) {gpu.launch kernel};
    \node[box] (t4) at (10.6,-0.9) {JIT + run\\\normalfont\itshape T4 GPU};
    \draw[arr] (op) -- node[lbl,above left=-2pt and -8pt]
      {convert-blas-to-scf\\Part 1} (scf);
    \draw[arr] (scf) -- node[lbl,above]{7 passes} (cpu);
    \draw[arr,draw=scaccent] (op) -- node[lbl,below left=-2pt and -8pt]
      {convert-blas-to-gpu\\this part} (gpu);
    \draw[arr] (gpu) -- node[lbl,below]{gpu-lower-to-nvvm-pipeline} (t4);
  \end{tikzpicture}
  \end{center}
  \vspace{0.3em}
  \begin{itemize}
    \item<2-> One op definition, \alert{two lowerings}: loops on the CPU in
      Part 1, a kernel on the GPU in this part.
    \item<3-> The pattern decides how the work maps onto the GPU: one thread
      per output element, x on the columns.
    \item<4-> It builds the kernel you read in step 1, and that runs through
      the same GPU pipeline.
  \end{itemize}
\end{frame}
\note{End here.}

\end{document}
